\section{Описание}

\subsection{Источник данных}
Робот обкачивает русские Викиновости (\texttt{ru.wikinews.org}) --- новостной
проект на движке MediaWiki, того же, что и Википедия из первой лабораторной
работы. В отличие от первой работы, статьи здесь не выбираются случайно через
API, а обходятся по ссылкам: робот загружает обычные HTML-страницы
\texttt{/wiki/<название>} в том виде, в каком их получает браузер
(\enquote{сырой} HTML со всей вёрсткой). Авторизация не нужна. Запросы
отправляются с идентифицирующим \texttt{User-Agent} (с адресом для связи), как
требуют правила Викимедиа, файл \texttt{robots.txt} сайта соблюдается.

\subsection{Общее устройство}
Робот --- программа \texttt{main.py} на Python 3 (библиотеки \texttt{requests},
\texttt{BeautifulSoup}, \texttt{pymongo}, \texttt{PyYAML}); вся логика --- один
класс \texttt{Crawler}. Единственный аргумент программы --- путь до YAML-конфига.
Хранилище --- MongoDB, две коллекции: \texttt{documents} (обкачанные документы)
и \texttt{documents\_queue} (очередь URL).

Алгоритм --- обход в ширину:
\begin{enumerate}
    \item в очередь добавляются стартовые URL (главная страница);
    \item из очереди берётся самый старый необработанный URL, страница
    загружается;
    \item документ сохраняется в \texttt{documents};
    \item из тела статьи извлекаются ссылки; подходящие URL, которых ещё нет в
    очереди, добавляются в неё;
    \item URL помечается обработанным, выдерживается пауза, шаг 2 повторяется,
    пока очередь не опустеет или не будет достигнут лимит документов.
\end{enumerate}

\subsection{Конфигурация}
В конфиге три обязательные секции: \texttt{db} (подключение к MongoDB),
\texttt{logic} (параметры робота) и \texttt{source} (источник: название, домен,
стартовые URL, закрытые пространства имён); необязательная секция
\texttt{logging}.

\begin{lstlisting}
db:
  host: localhost
  port: 27017
  database: wikinews
  collection: documents        # queue lives in "<collection>_queue"

logic:
  user_agent: "info-search-lab2/1.0 (student crawler; <e-mail>)"
  delay_between_requests: 1.0  # seconds between requests
  max_documents: 30000         # 0 = no limit
  refresh_interval_days: 7     # re-check documents older than this
  retry_attempts: 3            # HTTP-level retries and attempts per URL
  timeout_seconds: 30

source:
  name: Wikinews
  domain: ru.wikinews.org
  seeds: [<main page URL>]
  skip_namespaces: [<Category, Special, Talk, User, File, ...>]

logging: {level: INFO, file: crawler.log}
\end{lstlisting}

\subsection{Схема данных в MongoDB}
Коллекция \texttt{documents}, один документ на страницу:

\begin{tabular}{|p{4.6cm}|p{10.2cm}|}
\hline
Поле & Описание \\ \hline
\texttt{url} & нормализованный канонический URL; уникальный индекс \\ \hline
\texttt{html\_content} & \enquote{сырой} HTML страницы \\ \hline
\texttt{source} & название источника (\texttt{Wikinews}) \\ \hline
\texttt{crawl\_timestamp} & дата обкачки, Unix time stamp \\ \hline
\texttt{content\_hash} & MD5 текста статьи (для проверки изменений) \\ \hline
\texttt{title}, \texttt{content\_length}, \texttt{text\_length} & заголовок,
размер сырого HTML в байтах, длина текста статьи в символах \\ \hline
\texttt{status\_code}, \texttt{last\_modified} & код ответа и заголовок
\texttt{Last-Modified} (нужен для условного запроса при переобкачке) \\ \hline
\texttt{history} & предыдущие версии изменившихся документов (хэш и дата) \\ \hline
\end{tabular}

Коллекция \texttt{documents\_queue}, одна запись на URL: \texttt{\_id} --- URL,
\texttt{state} (\texttt{new} / \texttt{done} / \texttt{failed} / \texttt{error}),
\texttt{tries} (число неудачных попыток), \texttt{refresh} (URL взят на
переобкачку), \texttt{added} и \texttt{done\_at} (время постановки в очередь и
обработки).

\subsection{Очередь URL и защита от повторных загрузок}
\begin{itemize}
    \item \textbf{Нормализация.} Схема приводится к \texttt{https}, хост ---
    к нижнему регистру, отбрасываются фрагмент (\texttt{\#...}) и query,
    пробелы заменяются на \texttt{\_}, процентное кодирование приводится к
    единому виду. Ссылки \texttt{/wiki/A\%20B} и \texttt{/wiki/A\_B} ---
    один и тот же URL.
    \item \textbf{Очередь хранится в MongoDB}, \texttt{\_id} записи --- сам
    URL, добавление идёт через \texttt{upsert} с \texttt{\$setOnInsert}.
    Поэтому URL попадает в очередь один раз за всё время жизни базы, а не
    только за один запуск. Для документов дополнительно создан уникальный индекс
    по \texttt{url}.
    \item \textbf{Фильтр ссылок.} Берутся только ссылки из тела статьи
    (\texttt{\#mw-content-text}; меню и подвал сайта не обходятся) на домен
    \texttt{ru.wikinews.org} с путём \texttt{/wiki/}. Служебные пространства
    имён (Категория, Служебная, Обсуждение, Участник, Файл, Шаблон и т.\,д.)
    исключены. Просто отбросить все URL с двоеточием нельзя: заголовки
    новостей часто содержат двоеточие (<<Александр Егоров: Нам не пришлось
    что-то выдумывать>>), поэтому сравнивается именно префикс до двоеточия со
    списком пространств имён из конфига.
    \item \textbf{robots.txt} загружается один раз при старте (нашим
    \texttt{User-Agent}) и проверяется при добавлении URL в очередь.
    \item \textbf{Канонические URL.} Если страница открылась по алиасу или
    перенаправлению, документ сохраняется под URL из
    \texttt{<link rel="canonical">}, а этот канонический URL помечается в очереди
    обработанным, чтобы страница не загружалась дважды.
    \item \textbf{Порядок обхода} --- FIFO по полю \texttt{added}; URL, на
    которых были сбои, уходят в конец очереди (сортировка по \texttt{tries}).
\end{itemize}

\subsection{Обработка ошибок HTTP}
Ни одна ошибка одной страницы не прерывает обкачку:

\begin{tabular}{|p{5.2cm}|p{10cm}|}
\hline
Ситуация & Реакция \\ \hline
200, \texttt{text/html} & документ сохраняется \\ \hline
304 Not Modified & документ не изменился, обновляется только
\texttt{crawl\_timestamp} \\ \hline
404, 410 и другие 4xx & URL помечается \texttt{failed} и больше не
запрашивается \\ \hline
200, но не HTML & \texttt{failed} \\ \hline
429, 5xx & библиотека \texttt{urllib3} повторяет запрос
\texttt{retry\_attempts} раз с нарастающей паузой (1, 2, 4~с), заголовок
\texttt{Retry-After} учитывается \\ \hline
таймаут, обрыв соединения, 429/5xx после всех повторов & \texttt{tries}
увеличивается, URL уходит в конец очереди; после \texttt{retry\_attempts}
неудач --- состояние \texttt{error}, при следующем запуске такие URL
возвращаются в очередь \\ \hline
\end{tabular}

\subsection{Возобновление работы}
Контрольная точка --- сама очередь. URL получает состояние \texttt{done} только
после того, как документ записан в базу. Поэтому робот можно остановить в любой
момент: \texttt{Ctrl+C} или \texttt{SIGTERM} (робот дозагружает текущую
страницу и выходит; повторный сигнал завершает его сразу), а при аварийном
завершении (даже \texttt{kill -9}) теряется максимум одна страница --- она
останется в состоянии \texttt{new} и будет загружена заново. При повторном
запуске робот берёт из очереди необработанные URL и продолжает с того места, где
остановился; стартовые URL добавляются идемпотентно.

\subsection{Переобкачка изменившихся документов}
При каждом запуске URL в состоянии \texttt{done}, обкачанные более
\texttt{refresh\_interval\_days} дней назад, возвращаются в очередь с признаком
\texttt{refresh}. Запуск робота по расписанию (например, \texttt{cron}) даёт
периодическую переобкачку. Для таких страниц отправляется условный запрос
\texttt{If-Modified-Since} (значение берётся из сохранённого
\texttt{Last-Modified}):
\begin{itemize}
    \item 304 --- обновляется только \texttt{crawl\_timestamp};
    \item 200 --- сравнивается MD5-хэш \emph{текста статьи}. Хэш считается не
    от всего HTML, поскольку в странице MediaWiki есть служебные данные
    (отчёт парсера, время генерации), которые могут меняться без изменения
    статьи. Совпал --- обновляется только дата; не совпал --- документ
    перезаписывается, а прежние хэш и дата добавляются в \texttt{history}.
\end{itemize}
Когда число документов достигло \texttt{max\_documents}, робот берёт из очереди
только URL на переобкачку и новые страницы не загружает.

\subsection{Проверка}
Файл \texttt{test\_main.py} содержит пять автоматических проверок на
подставной HTTP-сессии и MongoDB в памяти (\texttt{mongomock}): обход с
дедупликацией и фильтрацией ссылок (категории, внешние домены, query, заголовок
с двоеточием), перенаправление в закрытое пространство имён (документ не сохраняется); остановка посреди обкачки с последующим продолжением (каждый URL
загружается ровно один раз), переобкачка (изменившийся документ получает
запись в \texttt{history}, неизменившийся --- нет), ошибки HTTP (404,
503, обрыв соединения). Кроме того, выполнен контрольный прогон по живому сайту
\texttt{ru.wikinews.org} в три этапа: первичная обкачка до 100 документов;
<<перезапуск>> с продолжением из очереди до 200 документов; переобкачка всех
документов, у которых срок давности принудительно считался истёкшим.

\subsection{Статистика контрольного прогона}
Значения получены командой \texttt{python3 main.py config.yaml -{}-stats} и по
логу прогона; задержка между запросами 1~с.

\begin{tabular}{|p{6.4cm}|p{8.4cm}|}
\hline
Показатель & Значение \\ \hline
Количество документов & 200 \\ \hline
Размер \enquote{сырых} данных & 25{,}1 МБ \\ \hline
Размер выделенного текста & 1{,}5 МБ \\ \hline
Средний размер документа & 125{,}7 КБ \\ \hline
Средний объём текста в документе & 7{,}5 КБ \\ \hline
Скорость обкачки & 100 док.\ за 119~с и 100 док.\ за 139~с
($\approx$0{,}78 док/с, $\approx$1{,}3~с на документ) \\ \hline
Ошибки HTTP и сбои соединения & 0 \\ \hline
URL в очереди после 200 документов & 222 обработано, 3497 ожидают \\ \hline
Из них адреса-алиасы (канонический URL другой) & 22 из 222 \\ \hline
Повторная загрузка при перезапуске & 0 (каждый URL запрошен один раз) \\ \hline
Переобкачка 200 документов & 178 ответов 304, 22 ответа 200 с тем же
хэшем, изменившихся 0 \\ \hline
\end{tabular}

Очередь после 200 документов насчитывает уже около 3700 адресов: обход по
ссылкам из тела статьи быстро выходит за пределы стартовой страницы, так что
лимит в 30000 документов достижим. Страницы-алиасы --- в основном ссылки на
перенаправления: например, заголовок с запятой и без неё ведёт на один и тот же
документ; в таких случаях сохраняется ровно один документ под каноническим
URL. Если канонический URL указывает в закрытое пространство имён (например,
перенаправление в категорию), документ не сохраняется.

Для всего корпуса статистика получается той же командой по базе после полной
обкачки:

\begin{center}\texttt{python3 main.py config.yaml -{}-stats}\end{center}

\pagebreak
